\documentclass{article}
\usepackage{amsmath}

% Size and space
\usepackage{titlesec}
\titleformat{\subsection}{\small\bfseries}{\thesubsection}{1em}{}
\titlespacing*{\subsection}{0pt}{24pt}{12pt}

\begin{document}

Im folgenden gehen wir von einem 8-Bit Zweierkomplement-System aus.

\subsection*{1. Stellen Sie die Zahlen $36_{10}$, $-36_{10}$ und $-128_{10}$ dar. Inwiefern ändert sich die Zahldarstellung, wenn man mehr als 8 Bits verwenden würde?}

\begin{center}
\renewcommand{\arraystretch}{1.5}
\begin{tabular}{|c|c|c|}
\hline
Dezimal & 8-Bit & 10-Bit \\
\hline
$36_{10}$ & 00100100 & 0000100100 \\
$-36_{10}$ & 11011100 & 1111011100 \\
$-128_{10}$ & 10000000 & 1110000000 \\
\hline
\end{tabular}
\end{center}

Bei mehr Bits werden negative Zahlen von links mit Einsen aufgefüllt.

\subsection*{2. Bestimmen Sie die kleinste und größte darstellbare Zahl.}

$$
MIN_8 = -128_{10} = 10000000
$$
$$
MAX_8 = 127_{10} = 01111111
$$

\subsection*{3. Erklären Sie, warum das Negieren der kleinsten Zahl nicht darstellbar ist.}

Im Zweierkomplement-System ist Negation definiert als $NOT(a) + 1$. Wenn wir das mit der kleinsten Zahl $-128$ machen, kommen wir wieder bei $-128$ raus!

$$
NEG(10000000) = NOT(10000000) + 1 = 01111111 + 1 = 10000000
$$

\subsection*{4. Zeigen Sie, dass für alle Zahlen $a$ im Zweierkomplement gilt: $-a = NOT(a) + 1$.}

Wenn wir jedes Bit einer Zahl $a$ negieren, erhalten wir eine Zahl $NOT(a)$, die zu $a$ addiert eine Zahl ergibt, die nur aus Einsen besteht: $a + NOT(a) = 11\ldots1$.

Wir wissen, dass im Zweierkomplement eine Zahl nur aus Einsen immer den Wert $-1$ hat. Also sehen wir dass
$$
a + NOT(a) = -1 \iff -a = NOT(a) + 1
$$

\newpage

Folgend gehen wir von IEEE-754 single precision aus.

\subsection*{1. Stellen Sie die IEEE-754 Zahl $C27FC000$ als Dezimalzahl dar.}

Zunächst stellen wir die Zahl in binärer Form dar. Da das Hexadezimalsystem $16 = 2^4$ als Base hat,
können wir jede Hex-Ziffer als eine Folge von 4 Bits darstellen. Wenn wir diese aneinanderreihen, erhalten wir die 
ganze Zahl. Mit $C = 12 = 1100$, $2 = 0010$, $7 = 0111$ und $F = 15 = 1111$ ergibt das die folgende 32-bit Zahl:

$$11000010011111111100000000000000$$

Aufgeteilt in Vorzeichen-Bit($s$), Mantisse($m$) und Exponent($e$):
\begin{itemize}
    \item $s = 1$
    \item $e = 10000100 - 127 = (128 + 4) - 127 = 5$
    \item $m = 1 + 0.11111111100000000000000_2 = 1 + \frac{511}{512} =  1.99804687$
\end{itemize}

Damit berechnen wir unsere Zahl laut Formel $(-1)^s \cdot m \cdot 2^e$:

$$(-1) \cdot 1.99804687 \cdot 2^5 = -63.9375$$

\subsection*{2. Stellen Sie die Fließkommazahl $42.875$ als Binärzahl dar.}

Wir brechen zunächst auf in ganze Zahl und Bruch:
$$42 = 32 + 8 + 2 = 101010_2$$
$$0.875 = \frac{7}{8} = 4/8 + 2/8 + 1/8 = 1/2 + 1/4 + 1/8 = 0.111_2$$

Zusammen ergibt das in normalisierter Form: 
$$42.875 = 101010.111_2 = 1.01010111_2 \cdot 2^5$$.

Wir erinnern uns, dass single precision Zahlen repräsentiert werden als
$$(-1)^s \cdot (1+\text{Mantisse}) \cdot 2^{\text{Exponent} - 127}$$ und kommen in unserem Fall auf
$$(-1)^0 \cdot (1 + 0.01010111000000000000000) \cdot 2^{132-127}$$.

In single precision repräsentieren wir die Zahl $42.875$ also durch
$$ \underbrace{0}_{\text{Vorzeichen}} \; \underbrace{10000100}_{\text{Exponent}} \; \underbrace{01010111000000000000000}_{\text{Mantisse}} $$


$$42 = 101010_2, \quad \frac{7}{8} = 0.111_2$$
$$42\frac{7}{8} = 101010.111_2$$
\end{document}